\chapter*{Încadrarea în literatura de specialitate}
\addcontentsline{toc}{chapter}{Încadrarea în literatura de specialitate}

\section*{Recunoașterea emoțiilor faciale folosind rețele SNN bio-inspirate}

Lucrarea elaborată de Mansouri-Benssassi și Ye \cite{mansouri2021} investighează aplicabilitatea rețelelor neuronale cu impulsuri (SNN) bio-inspirate în sarcini de recunoaștere a emoțiilor faciale (FER), punând un accent deosebit pe robustețea la zgomot și capacitatea de generalizare. Arhitectura propusă adaptează modelul Diehl și Cook \cite{diehl2015}, utilizând învățarea nesupervizată prin plasticitate dependentă de timpul impulsurilor (STDP). Procesarea imaginilor implică utilizarea filtrului Laplacian of Gaussian (LoG) pentru extracția eficientă a marginilor, urmată de o codificare de tip \textit{rate coding} prin trenuri de impulsuri Poisson, unde intensitatea pixelilor dictează rata de declanșare a neuronilor de tip Leaky Integrate-and-Fire (LIF). Modelul este structurat pe trei straturi: unul de intrare, un strat excitator (cu funcții de convoluție) și un strat inhibitor pentru selecția neuronului activ. Evaluările realizate pe seturile de date JAFFE și CK+ relevă faptul că SNN-ul obține performanțe comparabile cu metodele HOG+SVM și superioare arhitecturilor CNN pe seturi de date mici. Un aspect critic evidențiat este superioritatea SNN în experimentele de generalizare \textit{cross-dataset} (atingând până la 92\%) și reziliența remarcabilă la zgomot de tip \textit{salt and pepper}; la un nivel de degradare de 0,5, SNN menține o acuratețe de 56,2\%, depășind net modelele CNN (22,6\%) și SVM (14,1\%). Totuși, se identifică limitări privind utilizarea codificării pe bază de rată în detrimentul celei temporale și necesitatea rulării pe hardware neuromorfic dedicat pentru a maximiza eficiența energetică și computațională față de implementările actuale pe CPU/GPU.

\section*{Abordări bazate pe evenimente: Arhitectura Spiking-FER}

O contribuție majoră în domeniu este adusă de Barchid et al. \cite{barchid2023}, care introduc conceptul de „Event-based FER” și propun primul benchmark de acest tip prin conversia seturilor de date video clasice (CK+, ADFES, CASIA și MMI) în fluxuri de evenimente folosind simulatorul V2E \cite{yuhuang2021}. Autorii propun \textit{Spiking-FER}, o arhitectură SNN convoluțională profundă bazată pe SEW-ResNet-18, care utilizează neuroni de tip \textit{Integrate-and-Fire} (IF). Pentru a permite antrenarea de tip \textit{end-to-end} și a depăși problema non-diferențiabilității funcției de prag, se utilizează \textit{Surrogate Gradient Learning} cu o aproximare bazată pe funcția arctan. Spre deosebire de metodele modulare, abordarea \textit{end-to-end} permite optimizarea simultană a extractorului de trăsături și a clasificatorului. Rezultatele experimentale indică faptul că \textit{Spiking-FER} atinge performanțe comparabile cu modelele ANN, însă cu o eficiență energetică net superioară, fiind de până la 65,39 ori mai eficient din punct de vedere computațional. Studiul analizează, de asemenea, impactul tehnicilor de augmentare a datelor specifice evenimentelor (\textit{Event Data Augmentation}), concluzionând că metodele \textit{EventDrop} și \textit{Mirror} sunt esențiale pentru reglarea antrenării pe seturi de date limitate. Totuși, o limitare importantă a studiului rămâne dependența de date simulate în locul senzorilor DVS (\textit{Dynamic Vision Sensor}) reali și utilizarea bazelor de date de laborator, care pot să nu reflecte complexitatea emoțiilor spontane din mediul real.

\section*{Fuziunea multimodală și arhitecturi evolutive: Modelul FusionSense}

Lucrarea prezentată de Tan et al. \cite{tan2020} introduce \textit{FusionSense}, primul sistem de recunoaștere a emoțiilor bazat pe rețele neuronale cu impulsuri care utilizează fuziunea datelor multimodale (expresii faciale corelate cu semnale fiziologice precum ECG, conductanța pielii și respirația). Autorii utilizează cadrul \textit{NeuCube} \cite{kasabov2014}, o arhitectură SNN evolutivă (eSNN) care încorporează o structură spațială 3D de tip \textit{reservoir} (3D-SNNr) cu conectivitate \textit{small-world}. Pipeline-ul vizual utilizează algoritmul Viola-Jones pentru detecție și metoda Kanade-Lucas-Tomasi (KLT) pentru urmărirea punctelor de interes, extrăgând distanțe relative între 68 de repere faciale (\textit{landmarks}) raportate la o stare neutră de referință. Un aspect metodologic distinct este utilizarea codificării de tip \textit{population coding} bazată pe câmpuri receptive Gaussiene, urmată de o învățare nesupervizată prin STDP și o etapă supervizată ce folosește reguli de tip \textit{Rank Order} (deSNN). Experimentele pe setul de date MAHNOB-HCI demonstrează că fuziunea la nivel de trăsături (\textit{feature-level fusion}) este superioară celei la nivel de decizie, atingând o acuratețe de 73,15\% pentru valența binară. Sistemul se remarcă prin capacitatea de antrenare incrementală (\textit{one-pass}), fiind optim pentru aplicații \textit{real-time}. Totuși, studiul prezintă limitări legate de excluderea semnalelor EEG și de dependența de o fereastră temporală fixă de 30 de secunde, care poate omite dinamica emoțională pe termen lung.

\section*{Recunoașterea identității prin dinamica mișcărilor \\ faciale}

Lucrarea propusă de Liu et al. \cite{liu2020} introduce o paradigmă inovatoare de re\-cu\-noaș\-te\-re facială bazată pe dinamica mișcărilor musculare, metodă eficientă chiar și în scenarii cu fețe deghizate sau parțial ocultate. Spre deosebire de sistemele convenționale care se bazează pe trăsături spațiale statice, acest model utilizează cadre diferențiale pentru a captura micro-mișcările unice induse de vorbire. Arhitectura este de tip SNN profund, structurată pe cinci straturi care simulează calea ventrală a sistemului vizual biologic, incluzând module de extracție a trăsăturilor de nivel înalt prin filtre Gabor și codificare temporală bazată pe \textit{Rank Order Coding} (ROC). Un element central al modelului este stratul de învățare care utilizează mecanismul \textit{Soft Winner-Take-All} (WTA) cu inhibiție laterală și reguli STDP \textit{event-driven}, asigurând selectivitatea neuronilor pentru tipare motorii specifice. Evaluările pe setul de date MakeFaceDB au demonstrat rate de succes între 95\% și 100\%, depășind modelele CNN echivalente, în special în scenariile cu deghizări complexe (măști, barbă falsă). Autorii subliniază eficiența procesului de convergență, care necesită doar două clipuri video pentru antrenare. Totuși, sunt identificate limitări practice importante: dependența de mișcarea activă a subiectului (necesitatea de a vorbi), sensibilitatea ridicată la mișcările de fundal sau ale camerei și necesitatea unei reglări manuale riguroase a hiperparametrilor pentru a menține stabilitatea rețelei.

\section*{Recunoașterea emoțiilor prin senzori de viziune dinamică simulați}

Lucrarea propusă de Wang et al. \cite{wang2022} abordează deficitul de seturi de date pentru senzori de viziune dinamică (DVS) în recunoașterea emoțiilor prin propunerea unei metode de generare a datelor simulate din segmente video convenționale (setul CMU-MOSI \cite{zadeh2016}). Autorii introduc un cadru SNN bazat pe modelul de neuron \textit{Parametric Leaky Integrate-and-Fire} (PLIF), care permite optimizarea automată a parametrilor de scurgere (\textit{leak}) în timpul antrenării. O contribuție metodologică importantă este utilizarea unei prime rame de tip \textit{float-value} pentru a stabiliza potențialul de membrană, urmată de trenuri de impulsuri binare, strategie care accelerează scăderea funcției de pierdere și îmbunătățește acuratețea față de modelele care folosesc exclusiv \textit{spikes}. Arhitectura este compusă dintr-un modul de extracție a trăsăturilor cu straturi convoluționale și \textit{average pooling}, un grup de neuroni de votare și un modul de mapare emoțională care traduce ratele de declanșare în etichete de polaritate. Antrenarea se realizează prin gradient surogat și \textit{backpropagation through time} (BPTT). Rezultatele experimentale demonstrează că modelul SNN reduce consumul de energie cu 75\% față de echivalentul ANN, atingând performanțe optime la un prag de tensiune de 0,4. Totuși, datele simulate nu pot reproduce fidel rezoluția temporală fină a unui senzor DVS real, iar constrângerea ferestrei temporale la 68 de cadre poate duce la pierderea informațiilor în cazul secvențelor video de lungimi inegale.